% =====================================================================
%  Chapitre 1 — Étude et analyse des besoins
% =====================================================================
\chapter{Étude et analyse des besoins}
\label{chap:analyse}

\section{Contexte du projet}

Le projet se déroule chez \varEntreprise, laboratoire pharmaceutique dont le jeu de
données de démonstration reflète l'organisation : quatre implantations — un site de
production à Tunis, un centre de R\&D à Lyon, une plateforme de distribution à
Casablanca et un siège à Paris — et huit départements, du Commercial aux Affaires
réglementaires.

Cette dispersion géographique et fonctionnelle est précisément ce qui rend
l'arbitrage difficile. Un besoin de prévision des ventes exprimé par le Commercial de
Tunis et un besoin de prévision de la demande formulé par la Supply Chain de
Casablanca reposent sur le même référentiel produit, mais rien ne le signale tant que
les deux dossiers vivent dans des tableurs distincts.

\section{Présentation de \appname}

\appname{} est une plateforme de pilotage de portefeuille, destinée à la direction, aux
responsables de département et aux fonctions transverses d'un laboratoire. Elle n'est
pas un outil de développement de modèles : elle intervient \emph{avant} le projet
technique, pour décider s'il doit exister, et \emph{après} sa mise en production, pour
vérifier qu'il tient ses promesses.

La plateforme est multi-entreprises. Chaque organisation cliente dispose de ses sites,
départements, utilisateurs, rôles et données métier, sans aucune visibilité sur les
autres.

\section{Problématique}

La question à laquelle le projet répond peut s'énoncer ainsi :

\begin{remarque}[Problématique]
Comment doter une organisation pharmaceutique d'un cadre unique et traçable qui
transforme des idées d'intelligence artificielle dispersées en un portefeuille de
projets arbitrés, gouvernés et dont la valeur est mesurée, tout en respectant les
exigences d'explicabilité et de conformité propres au secteur ?
\end{remarque}

Trois sous-questions en découlent :

\begin{enumerate}
  \item Comment qualifier objectivement le type de solution adapté à un besoin, sans
    céder à l'effet de mode qui pousse à répondre « IA » à toute question ?
  \item Comment comparer des opportunités issues de départements différents sur une
    grille commune, tout en laissant les pondérations ajustables selon la stratégie du
    moment ?
  \item Comment garantir qu'une décision prise il y a six mois reste justifiable devant
    un comité ou un inspecteur ?
\end{enumerate}

\section{Objectifs généraux}

\begin{enumerate}
  \item Centraliser le diagnostic de maturité, l'expression des besoins et le
    portefeuille de projets dans un référentiel unique.
  \item Produire des décisions explicables : toute note, tout score et toute
    recommandation doivent pouvoir être décomposés.
  \item Garantir l'étanchéité entre organisations clientes et le cloisonnement des
    droits à l'intérieur de chacune.
  \item Mesurer la valeur réellement captée et la confronter à la valeur promise.
\end{enumerate}

\section{Objectifs spécifiques}

\begin{itemize}
  \item Évaluer la maturité sur deux axes, technique et métier, avec pondération et
    niveau de confiance par dimension.
  \item Identifier les dimensions bloquantes, c'est-à-dire celles qui interdisent un
    passage en production tant qu'elles ne sont pas remédiées.
  \item Qualifier chaque opportunité parmi huit types de solution, avec justification
    écrite, alternatives écartées et garde-fous.
  \item Calculer un score de priorité sur six critères pondérables, diminué d'une
    pénalité de risque.
  \item Détecter les doublons et les synergies entre opportunités de départements
    différents.
  \item Générer une fiche projet versionnée et soumise à décision de comité.
  \item Suivre le cycle de vie par phases avec gates de passage.
  \item Calculer TCO, gains annuels, ROI et délai de retour, au projet et au
    portefeuille.
  \item Journaliser toute modification significative de façon immuable.
\end{itemize}

\section{Besoins fonctionnels}

Les besoins fonctionnels se regroupent en douze modules, notés M1 à M12. Le
tableau~\ref{tab:modules} en donne l'état réel dans l'application livrée. Cet état a
été établi en inspectant le code : routers FastAPI, services métier, écrans React et
suite de tests.

\begin{longtable}{@{}L{1.1cm} L{6.6cm} L{2.6cm} L{4.2cm}@{}}
\caption{État d'implémentation des douze modules fonctionnels}
\label{tab:modules}\\
\toprule
\textbf{Mod.} & \textbf{Description} & \textbf{État} & \textbf{Constat dans le code} \\
\midrule
\endfirsthead
\multicolumn{4}{c}{\small\itshape suite du tableau~\ref{tab:modules}}\\
\toprule
\textbf{Mod.} & \textbf{Description} & \textbf{État} & \textbf{Constat dans le code} \\
\midrule
\endhead
\bottomrule
\endfoot

M1 & Organisation, sites, départements, utilisateurs, rôles et permissions
   & \fait
   & 18 endpoints sous \code{/org}, 41 permissions, 11 rôles, écran Administration à
     trois onglets \\

M2 & Diagnostic AS-IS inspiré de TOGAF : campagnes et réponses pondérées
   & \fait
   & 8 endpoints \code{/diagnostic}, questionnaire éditable dans l'interface \\

M3 & Expression des besoins : questionnaire commun et volet départemental
   & \fait
   & \code{services/questionnaire.py}, 5 sections communes et 7 volets métier \\

M4 & Qualification du type de solution
   & \fait
   & \code{services/recommendation.py}, 8 types, 39 règles pondérées \\

M5 & Readiness : score sur 100, dimensions bloquantes
   & \fait
   & \code{services/readiness.py}, distinction bloquante / vigilance / à instruire \\

M6 & Croisements interservices : processus, applications, données, dépendances,
     doublons
   & \fait
   & 14 endpoints \code{/crossings}, écran Architecture \& données \\

M7 & Priorisation du portefeuille, pondérations paramétrables
   & \fait
   & \code{services/scoring.py}, panneau de repriorisation avec curseurs \\

M8 & Fiche projet : génération, versions, validation, export
   & \partiel
   & Génération, versionnement et décisions opérationnels ; export en JSON et HTML
     imprimable, pas de génération \code{.docx} \\

M9 & Suivi du cycle de vie : phases, jalons, tâches, risques, gates
   & \partiel
   & API complète (16 endpoints \code{/projects}) mais aucun écran de gestion de
     projet dédié \\

M10 & Pilotage de la valeur : TCO, gains, ROI, payback
   & \fait
   & \code{services/value.py}, 9 endpoints \code{/value}, écran Valeur \& ROI \\

M11 & Tableaux de bord et agrégations
   & \fait
   & 4 endpoints \code{/dashboards}, vues Direction et Département \\

M12 & Administration et journal d'audit immuable
   & \partiel
   & Journal complet et écran dédié ; notifications persistées mais sans écran de
     consultation \\
\end{longtable}

\begin{ecart}
Trois modules ne sont que partiellement couverts, et il serait malhonnête de les
présenter autrement.

\textbf{M8 — export.} Le cahier des charges prévoyait un export bureautique. Le bouton
de l'interface porte encore le libellé « Export Word / PDF » alors que l'API produit du
JSON et une page HTML mise en page pour l'impression. L'utilisateur obtient bien un PDF
en passant par la fonction d'impression de son navigateur, mais aucun fichier
\code{.docx} n'est généré. Le libellé est donc en avance sur la fonction.

\textbf{M9 — interface de suivi.} Les endpoints existent et sont protégés : phases,
jalons, tâches, risques, mitigations, incidents et passage de gates. Aucun écran dédié
ne les expose. Le PMO retrouve ses projets, leur avancement et leur prochaine gate dans
le tableau de bord de département, ce qui couvre le pilotage mais pas la gestion
détaillée des tâches.

\textbf{M12 — notifications.} La table et les endpoints existent, l'icône de cloche de
l'en-tête affiche une pastille de non-lu, mais aucun panneau ne permet de lire les
notifications.
\end{ecart}

\section{Besoins non fonctionnels}

\begin{longtable}{@{}L{3.4cm} L{10.8cm}@{}}
\caption{Besoins non fonctionnels et traitement retenu}
\label{tab:nfr}\\
\toprule
\textbf{Exigence} & \textbf{Traitement dans l'application} \\
\midrule
\endfirsthead
\toprule
\textbf{Exigence} & \textbf{Traitement dans l'application} \\
\midrule
\endhead
\bottomrule
\endfoot

Étanchéité des tenants
& Toute table métier porte \code{organization\_id}. Le jeton JWT transporte
  l'identifiant du tenant et la dépendance \code{get\_tenant\_id} le lit depuis le
  jeton, jamais depuis la requête. Neuf tests couvrent ce point. \\

Explicabilité
& Trois endpoints renvoient le raisonnement et non le seul résultat :
  \code{/opportunities/\{id\}/score}, \code{/opportunities/\{id\}/qualify} et
  \code{/diagnostic/readiness}. \\

Traçabilité
& Journal \code{audit\_log} en insertion seule : l'API n'expose ni mise à jour ni
  suppression sur cette table. Chaque entrée conserve auteur, horodatage, entité,
  ancienne et nouvelle valeur, et adresse IP. \\

Sécurité des accès
& Authentification par jeton JWT signé HS256, mots de passe hachés avec bcrypt,
  contrôle de permission sur chaque endpoint sensible. \\

Robustesse du schéma
& Migrations Alembic versionnées, écrites pour s'appliquer sur une base déjà peuplée
  sans perte de données. \\

Maintenabilité
& Logique métier isolée dans sept services sans dépendance à FastAPI, donc testables
  indépendamment du transport HTTP. \\

Ergonomie
& Interface en français, navigation et actions dérivées des permissions réelles,
  responsive jusqu'à 390~pixels de large. \\
\end{longtable}

\section{Contraintes}

\begin{itemize}
  \item \textbf{Contrainte réglementaire.} Le périmètre GxP impose que toute
    recommandation reste explicable et qu'une validation humaine subsiste. Le moteur de
    qualification ajoute automatiquement trois garde-fous lorsque l'opportunité est
    déclarée GxP : validation informatique, traçabilité ALCOA+ et revue humaine
    obligatoire.
  \item \textbf{Contrainte de confidentialité.} Deux laboratoires concurrents peuvent
    être hébergés sur la même instance. L'isolation ne peut pas dépendre du bon
    comportement du frontend.
  \item \textbf{Contrainte de langue.} L'interface et les données de démonstration sont
    en français, ce qui a imposé une attention particulière aux jeux de caractères,
    notamment pour l'export CSV destiné à Excel.
  \item \textbf{Contrainte d'environnement.} Le développement s'est fait sous Windows
    avec PowerShell, ce qui a influencé la documentation d'installation.
\end{itemize}

\section{Méthodologie adoptée}

Le développement a suivi une progression par couches, du socle vers l'interface, avec
validation à chaque étape.

\begin{enumerate}
  \item \textbf{Modélisation des données.} Les 34~tables ont été définies avant tout
    développement d'API, puis figées par une migration Alembic initiale.
  \item \textbf{Services métier.} Readiness, qualification, priorisation et
    valorisation ont été écrits comme des fonctions pures, testables sans base de
    données ni serveur.
  \item \textbf{Exposition HTTP.} Les routers FastAPI se contentent d'orchestrer :
    vérifier les droits, charger les entités, appeler le service, journaliser.
  \item \textbf{Jeu de démonstration.} Un seed réaliste a été construit tôt, pour que
    chaque écran soit développé face à des données crédibles plutôt qu'à des tableaux
    vides.
  \item \textbf{Interface.} Les écrans ont été développés un par un, chacun vérifié par
    capture dans un navigateur réel avant de passer au suivant.
  \item \textbf{Durcissement.} Le contrôle d'accès a été repris en profondeur dans un
    second temps, une fois les fonctionnalités stabilisées, et couvert par une suite de
    tests dédiée.
\end{enumerate}

Cette dernière étape mérite d'être soulignée : la première version du contrôle d'accès
ne comptait que dix permissions, toutes en écriture. Les quarante-six endpoints de
lecture n'étaient protégés que par l'authentification, ce qui laissait tout utilisateur
connecté lire l'intégralité de son organisation. La reprise a porté le catalogue à
41~permissions et introduit la notion de périmètre.
